% from ega2-1-fr.tex
\oldpage[II]{90}

l'assertion résulte alors de (\hyperref[II.4.5.2-fr]{4.5.2, b)}). Il est
clair que b) entraîne b')~; enfin, pour prouver que b') entraîne a), il
suffit de considérer un recouvrement de $Y$ par des ouverts affines
$U_\alpha$ et d'appliquer le critère
(\hyperref[II.4.5.2-fr]{4.5.2, b')}) à chaque faisceau
$\mathcal{L}|f^{-1}(U_\alpha)$.

Pour les deux dernières assertions, on utilise le fait que
$\sigma^\flat$ est ici l'identité, et l'explicitation des homomorphismes
(\hyperref[II.3.7.9.1-fr]{3.7.9.1}) et
(\hyperref[II.3.7.9.2-fr]{3.7.9.2})~; il en résulte aussitôt que
(\hyperref[II.4.6.3.1-fr]{4.6.3.1}) est un isomorphisme. Quant à
(\hyperref[II.4.6.3.2-fr]{4.6.3.2}), on peut se ramener au cas où $Y$ est
affine, donc $\mathcal{L}$ ample~; il est clair que l'homomorphisme
(\hyperref[II.4.6.3.2-fr]{4.6.3.2}) est injectif, et le critère
(\hyperref[II.4.5.2-fr]{4.5.2, c)}) montre qu'il est surjectif, d'où la
conclusion.

\begin{corollary}[4.6.4]
\label{II.4.6.4-fr}
Soit $(U_\alpha)$ un recouvrement ouvert de $Y$. Pour que $\mathcal{L}$
soit ample relativement à $Y$, il faut et il suffit que
$\mathcal{L}|f^{-1}(U_\alpha)$ soit ample relativement à $U_\alpha$ pour
tout $\alpha$.
\end{corollary}

La condition b) est en effet locale sur $Y$.

\begin{corollary}[4.6.5]
\label{II.4.6.5-fr}
Soit $\mathcal{K}$ un $\mathcal{O}_Y$-Module inversible. Pour que
$\mathcal{L}$ soit $Y$-ample, il faut et il suffit que
$\mathcal{L}\otimes f^*(\mathcal{K})$ le soit.
\end{corollary}

C'est une conséquence évidente de (\hyperref[II.4.6.4-fr]{4.6.4}) en
prenant les $U_\alpha$ tels que $\mathcal{K}|U_\alpha$ soit isomorphe à
$\mathcal{O}_Y|U_\alpha$ pour tout $\alpha$.

\begin{corollary}[4.6.6]
\label{II.4.6.6-fr}
Supposons $Y$ affine~; pour que $\mathcal{L}$ soit $Y$-ample, il faut et il
suffit que $\mathcal{L}$ soit ample.
\end{corollary}

C'est une conséquence immédiate de la définition
(\hyperref[II.4.6.1-fr]{4.6.1}), et des critères
(\hyperref[II.4.6.3-fr]{4.6.3, b)}) et
(\hyperref[II.4.5.2-fr]{4.5.2, b)}), car ici
$\operatorname{Proj}(\mathcal{S})
=\operatorname{Proj}(\Gamma(Y,\mathcal{S}))$ par définition.

\begin{corollary}[4.6.7]
\label{II.4.6.7-fr}
Soit $f:X\to Y$ un morphisme quasi-compact. Supposons qu'il existe un
$\mathcal{O}_Y$-Module quasi-cohérent $\mathcal{E}$ et un $Y$-morphisme
$g:X\to P=\mathbf{P}(\mathcal{E})$ qui soit un homéomorphisme de l'espace
sous-jacent à $X$ sur un sous-espace de $P$~; alors
$\mathcal{L}=g^*(\mathcal{O}_P(1))$ est $Y$-ample.
\end{corollary}

On peut en effet supposer $Y$ affine~; le corollaire résulte alors du
critère (\hyperref[II.4.5.2-fr]{4.5.2, a)}), de la formule
(\hyperref[II.3.7.3.1-fr]{3.7.3.1}) et de
(\hyperref[II.4.2.3-fr]{4.2.3}).

\begin{proposition}[4.6.8]
\label{II.4.6.8-fr}
Soient $X$ un schéma quasi-compact ou un préschéma dont l'espace
sous-jacent est noethérien, $f:X\to Y$ un morphisme séparé quasi-compact.
Pour qu'un $\mathcal{O}_X$-Module inversible $\mathcal{L}$ soit $f$-ample,
il faut et il suffit que l'une des conditions équivalentes suivantes soit
vérifiée~:
\begin{enumerate}
  \item[\rm{c)}] Pour tout $\mathcal{O}_X$-Module quasi-cohérent
  $\mathcal{F}$ de type fini, il existe un entier $n_0>0$ tel que, pour
  tout $n\geq n_0$, l'homomorphisme canonique
  $\sigma:f^*(f_*(\mathcal{F}\otimes\mathcal{L}^{\otimes n}))\to
  \mathcal{F}\otimes\mathcal{L}^{\otimes n}$ soit surjectif.
  \item[\rm{c')}] Pour tout faisceau d'idéaux quasi-cohérent de type fini
  $\mathcal{J}$ de $\mathcal{O}_X$, il existe un entier $n>0$ tel que
  l'homomorphisme canonique
  $\sigma:f^*(f_*(\mathcal{J}\otimes\mathcal{L}^{\otimes n}))\to
  \mathcal{J}\otimes\mathcal{L}^{\otimes n}$ soit surjectif.
\end{enumerate}
\end{proposition}

Comme $X$ est quasi-compact, il en est de même de $f(X)$ et il existe donc
un recouvrement fini $(U_i)$ de $f(X)$ par des ouverts affines de $Y$. Pour
démontrer la condition c) lorsque $\mathcal{L}$ est $f$-ample, on peut
remplacer $Y$ par les $U_i$ et $X$ par les $f^{-1}(U_i)$, car si on obtient
pour chaque $i$ un entier $n_i$ tel que c) ait lieu (pour $U_i$,
$f^{-1}(U_i)$ et $\mathcal{L}|f^{-1}(U_i)$) dès que $n\geq n_i$, il
suffira de prendre pour $n_0$ le plus grand des $n_i$ pour obtenir c) pour
$Y$, $X$ et $\mathcal{L}$. Or lorsque $Y$ est affine, la condition c)
découle de (\hyperref[II.4.5.5-fr]{4.5.5, d)}), compte tenu de
(\hyperref[II.4.6.6-fr]{4.6.6}). Il est trivial que c) entraîne c'). Enfin,
pour prouver que c') implique que $\mathcal{L}$ est $f$-ample, on peut
encore se borner au cas où $Y$ est affine~: en effet, tout faisceau
d'idéaux quasi-cohérent de type fini $\mathcal{J}_i$ de
$\mathcal{O}_X|f^{-1}(U_i)$ est la restriction d'un faisceau cohérent
d'idéaux de type fini de $\mathcal{O}_X$
(\hyperref[I.9.4.7-fr]{I, 9.4.7}) et l'hypothèse c') entraîne que

